\chapter{Environment Requirements}
\label{ch:envreqs}

\newthought{The \yamlkey{EnvReqs} block} is a safety gate. It says what a scan expects from the
machine---operating system, Python version, packages---and \Jarvis\ checks all of it \emph{before}
the run starts. A failed check stops the run early with a clear message, instead of crashing
halfway through.

\section{The minimal form}
\label{sec:env-minimal}

For most users this is enough:

\begin{yamlwin}
EnvReqs:
  OS:
    - name: linux
      version: ">=5.10.0"
    - name: Darwin
      version: ">=10.14"
  Python:
    version: ">=3.10"
    Dependencies: []
\end{yamlwin}

\noindent It says ``run on Linux $\geq 5.10$ or macOS (\code{Darwin}) $\geq 10.14$, with Python
$3.10$ or newer.'' That covers a typical setup.

\section{The keys}
\label{sec:env-keys}

\paragraph{\yamlkey{OS}} A list of acceptable operating systems. Each item has a \yamlkey{name}
(\code{linux} or \code{Darwin}) and a \yamlkey{version} condition. \Jarvis\ compares the current
machine against the list and stops if none match.

\paragraph{\yamlkey{Python}} The Python requirement. \yamlkey{version} is the minimum interpreter
version. \yamlkey{Dependencies} is a list of required packages, each with a \yamlkey{name},
\yamlkey{required} (\code{true}/\code{false}), and minimum \yamlkey{version}:

\begin{yamlwin}
  Python:
    version: ">=3.10"
    Dependencies:
      - name: numpy
        required: true
        version: ">=1.24"
      - name: scipy
        required: true
        version: ">=1.10"
\end{yamlwin}

\noindent If a required package is missing or too old, \Jarvis\ stops and prints the exact
\cli{pip install -U ...} command to fix it.

\paragraph{\yamlkey{Check\_default\_dependencies}} An optional way to reuse a shared environment
file. When \yamlkey{required: true}, \Jarvis\ loads the \textsc{yaml} at
\yamlkey{default\_yaml\_path} and merges its \yamlkey{EnvReqs} before validating:

\begin{yamlwin}
  Check_default_dependencies:
    required: true
    default_yaml_path: "&J/deps/environment_default.yaml"
\end{yamlwin}

\noindent New projects ship a \file{deps/environment\_default.yaml} for exactly this purpose
(Chapter~\ref{ch:quickstart}).

\paragraph{\yamlkey{CERN\_ROOT}} An optional block for scans that need CERN \textsc{root}. Add it
only if your tooling really depends on \textsc{root}; it can check the version and run a discovery
command:

\begin{yamlwin}
  CERN_ROOT:
    required: true
    version: ">=6.28"
    get_path_command: "root-config --prefix"
    Dependencies: []
\end{yamlwin}

\begin{jtip}
Start with just \yamlkey{OS} and \yamlkey{Python}. Add \yamlkey{Check\_default\_dependencies} when
your group maintains a shared default card, and \yamlkey{CERN\_ROOT} only when a scan truly needs
\textsc{root}. Less is more here---each requirement is one more thing that can block a run.
\end{jtip}

\section{A fuller example}
\label{sec:env-example}

\begin{yamlwin}
EnvReqs:
  OS:
    - name: linux
      version: ">=5.10.0"
    - name: Darwin
      version: ">=10.14"
  Python:
    version: ">=3.10"
    Dependencies:
      - name: numpy
        required: true
        version: ">=1.24"
      - name: xslha
        required: true
        version: ">=0.2"
  Check_default_dependencies:
    required: true
    default_yaml_path: "&J/deps/environment_default.yaml"
\end{yamlwin}

\section{Summary}
\label{sec:env-summary}

\yamlkey{EnvReqs} checks the platform before a run. Declare acceptable \yamlkey{OS} entries and a
\yamlkey{Python} version (plus any required packages); optionally merge a shared default card or
require \textsc{root}. It is the difference between a clean ``please upgrade NumPy'' message now
and a confusing crash an hour into a scan.
